\documentclass[book.tex]{subfiles}

\begin{document}

\chapter{Designing New Superconductors}

\label{chap:designing_superconductors}
\noindent\rule{11cm}{0.4pt} \\
\textbf{Prerequisites:}  \autoref{chap:schrodinger} \\
\textbf{Difficulty Level:} ** \\
\noindent\rule{11cm}{0.4pt}
\vspace{5mm}

The first superconductor was discovered a century ago in 1911. With a century's research, scientists have discovered some classes of higher temperature superconductors such as cuprates that reach a critical temperature at 133-138 K. Other superconductors exhibit superconductivity at higher pressures, with the current record holder lanthanum hydride at a critical temperature of 250 K at a pressure of 170 GPa. However, we remain far from the design of a room temperature, room pressure superconductor.

We have studied a variety of deep learning techniques which permit for the prediction of complex material properties starting from material structure. In this chapter we apply some of the techniques we've learned to do a case study of the problem of predicting the critical temperature of a prospective superconductor. We will start this chapter with a brief overview of the classical theory of superconductors, followed by a review of some recent research into the use of deep learning techniques to predict critical temperature of superconductors.


\section{Convolutional Featurizations}
Recent work \cite{konno2018deep} has proposed the use of deep learning architecture for superconducting materials design, trained it from a historical database of known superconductors. %They then test this model (trained on historical data) on a variety of recently discovered superconducting materials to see whether their model could have aided in materials discovery.

The base discovery rate for superconducting materials in screens is low, with a recent screen for superconducting materials finding that only 3\% of materials screened were superconducting \cite{konno2018deep}. These numbers match low activity rates in pharmaceutical screening assays. At present screening depends on experimental intuition and is mostly trial and error. %One of the broad dreams of scientific deep learning is to take human gut feeling analyses like these and turn them into reproducible processes. The design of a room temperature superconductor is one of science’s grand challenges, so it would be amazing to see deep learning play a role in solving this open problem. 

We have no training data yet that gives us strong clues on higher temperature superconductor theory, so we face a strong generalization challenge\cite{konno2018deep}. %This is an interesting bit of phrasing since zero-shot learning is more often used in tasks like translating a new language with no existing training data. 
Although classical physical theories like BCS theory provide some guidance, known higher temperature superconductors such as cuprates have complex strong electron correlations that foil classical computational techniques.
% (\textcolor{red}{TODO: Cite})

\subsection{Periodic Table Featurization}
For many potential superconducting materials, access to a spatial representation of the superconductors in question is not available, so featurization proves to be a channel. One idea is to simply encode materials as entries in the periodic table \cite{konno2018deep}. The basic idea is to take the periodic table as a subset of a 2 dimensional rectangle and then encode the fraction of the compound that belongs to each element. This representation is a form of two-dimensional one-hot encoding. We then do one additional split into s/p/d/f orbitals to represent the orbital characteristics of the valence electrons. This split allows the networks to learn to compute on the valence orbital blocks.

\begin{figure}
    \centering
    % \includegraphics{figures/Differentiable Physics/Materials ML/Designing Superconductors/periodic_table_featurization.png}
    \includegraphics{figures/Differentiable Physics/Materials ML/Designing Superconductors/Periodic table Feautization_SC (1).png}
    \caption{Periodic Table Featurization for Superconducting Materials. }
    \label{fig:periodic_featurization.}
\end{figure}
% \textcolor{red}{TODO: Replace}

This method is analogous to “imagification” strategies used elsewhere in scientific deep learning to process molecules \cite{goh2017chemception},  genetic sequences \cite{movva2019deciphering}, and microscopy \cite{christiansen2018silico} data. The periodic table was designed to structurally encode useful properties in a way that the human visual system could process. Given the relationships between convolutional neural networks and the human visual cortex, it’s not a far stretch to re-use this visual encoding to understand material structure.

\subsection{Dataset Augmentation}
%The network is trained on inorganic super conductor data from the SuperCon database (\textcolor{red}{TODO: Cite}) of known superconductors. There is processed version (\url{https://github.com/vstanev1/Supercon}) of SuperCon that was used in another paper \cite{stanev2018machine}, but which doesn’t contain time-split information (which we will see was used extensively in this paper). The preliminary model they trained achieves .92 test $R^2$ accuracy \cite{konno2018deep}.


%Models were trained with Adam for 6,000 epochs with batch size 32 and learning rate $2 \times 10^{-6}$. Models were 64 convolutional layers deep. Training took about 50 hours on their setup.
Just training on known superconductor structures will fail to calibrate models since there will be insufficient negatives. Earlier work found that models trained only on known superconductors have a high false positive rate. Unfortunately, existing superconductor datasets don't usually measure non-superconductors. 
%The paper reasons that this is likely due to the fact that the SuperCon dataset has very few non-superconductors (only 60), so the model has a bias from its training data.

%The paper applied this model to make predictions on the [crystallography open database](http://www.crystallography.net/cod/) (COD). The authors found that their trained models predicted that 17,000 compounds from COD were predicted to have $T_c$ (superconducting temperatures) over 10K, which is physically unrealistic. 

To solve this issue, we can use a data augmentation strategy. We can use compounds the from Crystallography Open Database that are not in the SuperCon superconductor dataset and reason that these compounds are unlikely to be superconducters, and so add them as training samples with $T_c = 0$. %This figure below summarizes the data sourcing strategy:

%\textcolor{red}{TODO ADD FIGURE}

%Here’s the joint training scheme rendered with this data augmentation strategy

%\textcolor{red}{TODO ADD FIGURE}

\section{Time Split Evaluation}

The model was evaluated on another dataset of materials discovered since 2010 . This dataset contains 400 compounds with 330 non-superconductors. To remove potential information leakage, the authors time-split SuperCon and COD and use only compounds added before 2010 as training data.% Here are some results from their models on this dataset

% TODO:CITE: ([reference](\url{https://iopscience.iop.org/article/10.1088/1468-6996/16/3/033503/meta}))

%\textcolor{red}{TODO ADD FIGURE}

The baseline was randomly selecting any compound from this dataset and predicting it to have $T_c > 0$. Comparing against a Random Forest baseline which uses manually designed features such as atomic mass, band gap, atomic configuration and melting temperature, the convolutional model achieves strong performance. However, the method does not generalize well to $T_c > 10$ predictions (since only a few superconductors in the training datat have this higher superconducting temperature). The baseline also does much worse in this case. The original work notes that precision and recall is much higher for the deep learned model over the baseline in both cases. %Comparison against the random forest is a little more muddled, but it’s not unreasonable to say the deep learning model has stronger performance (although I’d have liked to see error bars here).

%The authors also re-test the trained model against the COD database. I was a little confused by this comparison since COD was made part of the training data! I think the authors were looking to see what the trained model predicted on its own training data to see where the model predicted something that didn’t match the synthetic label of $T_c = 0$. 
%The original paper \cite{konno2018deep} also re-tests the trained model against the COD database. This test identifies $\textrm{CaBi}_2$ with this technique, a recently discovered superconductor not in SuperCon.

As a third test, the paper use training data sourced from before 2008 to see if they can find high-$T_c$ $\textrm{Fe}$ based superconductors (FeSCs), which were first discovered in 2008. Two precursor materials, LaFePO and LaFePFO were also removed from the training data since their discovery in 2006 spurred the development of FeSCs. 1,399 FeSCs known as of 2018 were used as the test set. 130 different training runs were used, which lead to stochastic variation in the predictions. The histogram below plots the number of predicted FeSCs on the test set from these various trained models. The work note that when they used a shallower 10 layer network (which had similar precision and $R^2$ as their models), they failed to identify FeSCs, and suggests that the depth of their networks may lead to greater generalizability.

%\textcolor{red}{TODO ADD FIGURE}

%The authors suggest that had their models been used earlier, they could have aided in the discovery of FeSCs, but acknowledge that the model spread noted in the figure above could have meant FeSCs were a lower priority target. The authors mention that they’re planning to attempt to use crystal structure data in future work to see if they can achieve higher predictive performance.

The use of time splits to validate the generalizability of the model is critical. %The variety of different test sets makes a stronger case that deep learning can serve as a useful tool in superconductor discovery. Here’s a table that summarizes the various train/test splits that the authors used to validate their models

%\textcolor{red}{TODO ADD FIGURE}

%The authors note that they plan to open source their code, but are waiting until they can run a few experiments with novel compounds that they’ve predicted. I hope that they also open source a clean form of their datasets. The raw SuperCon database doesn’t have a clean download option, and it might be a little tricky to reproduce their training/test splits without this additional information. It would also be very convenient to have nice versions of their FeSCs and other test datasets.

%\section{Generative Models}

\section{Limitations}

This method doesn't encode much of the known physics of superconductors, so there will be strong limitations on the generalizability to entirely unknown materials will be limited.


\end{document}